% 原书 p.42–p.45（PDF p.50–p.53）
% >>> p050
\section{可积系统（Integrable System）}

\begin{keytakeaways}
\begin{itemize}
  \item 物理图像：可积系统即「有足够多」相互对合的守恒量（首次积分）的哈密顿系统，
        其轨道被限制在不变环面上；Kepler 问题与环面作用是两大范例。
  \item 数学工具：Liouville 意义下的可积性（定义 2.6.1）、Poisson 交换的哈密顿函数组、
        角动量 $A=q\times p$ 与 Runge--Lenz 向量 $V_{\mathrm{RL}}$。
  \item 几何观点：$\{\dd H_i\}$ 在某开稠密集上线性无关，等价于哈密顿向量场
        $X_{H_i}$ 线性无关；环面上的矩映射分量给出 $n$ 个对合守恒量。
  \item 本节结论：Kepler 问题由 $H$、$A_3$ 与 $V_{\mathrm{RL}}$ 的分量构成可积系统；
        $\C^n$ 的 $\T^n$-作用与辛环面流形都自动给出可积系统（例 2.6.3），
        而平面三体问题一般不可积（例 2.6.4）。
\end{itemize}
\end{keytakeaways}

让我们用三元组 $(M,\omega,H)$ 表示一个哈密顿系统，其中 $(M,\omega)$ 是辛流形，$H$ 是其上
的哈密顿函数。本节我们将介绍该系统在 \emph{Liouville 意义下}的可积性；可积性（局部地）
等价于作用量–角坐标系统的存在。稍后我们将从几何角度看到（见定理 2.8.1），可积哈密顿
系统具有相当简单的结构。

\begin{wdefn}{2.6.1}
给定哈密顿系统 $(M^{2n},\omega,H)$，若存在 $H_1(=H),H_2,\cdots,H_n$，使得集合
$\{\dd H_i\}_{i=1}^n$ 在 $M$ 的某开稠密子集上线性无关，且 $H_1,H_2,\cdots,H_n$
满足如下 Poisson 交换条件
\[
  \{H_i,H_j\}=0\qquad\text{对任意 } 1\le i,j\le n,
\]
则称该系统是（完全）\textbf{可积的}（completely integrable）。
\end{wdefn}

\begin{wrem}{2.6.1}
定义 2.6.1 中的可积性定义相当强。微分 $\{\dd H_i\}_{i=1,\cdots,n}$ 线性无关（几乎处处）
意味着诱导出的哈密顿向量场 $\{X_{H_i}\}_{i=1,\cdots,n}$ 线性无关（几乎处处）。此外，
Poisson 交换条件意味着这些 $X_{H_i}$ 局部可构成一个标准/平坦的 $n$ 维坐标系：沿任一
方向运动都不影响其他方向。在某种程度上，这启发了定理 2.8.1 所提供的作用量–角坐标。
\end{wrem}
% [Tutor Note: 这里的「几乎处处」应理解为「在某个开稠密子集上」而非「除零测集外」——
% 定义 2.6.1 原文用 open dense subset，它与「流形上的稠密开集」一致。等价地：
% $\{\dd H_i\}$ 线性无关 $\iff$ 由 $\omega$ 对偶给出的 $X_{H_i}$ 线性无关，因为
% $v\mapsto\iota_v\omega$ 是 $T_pM\to T^*_pM$ 的线性同构（非退化性）。]

% >>> p051
可积系统的一个例子是空间 Kepler 问题，它是经典力学中二体问题的特例。在此情形下，两体
之间的相互作用由中心力 $F$ 支配，它随两体间距离的平方反比变化。

\begin{wexam}{2.6.1}
（空间 Kepler 问题）对由
\begin{equation}
  H(q,p)=\frac12|p|^2-\frac{1}{|q|}
  \tag{2.31}
\end{equation}
定义的哈密顿函数 $H:T^*(\R^3\setminus\{0\})\to\R$，在标准辛结构
$\omega_{\mathrm{can}}$ 下，由 (2.6) 中的哈密顿方程算得
\[
  \begin{cases}
    \dot q_i=\dfrac{\partial H}{\partial p_i}=p_i\\[2mm]
    \dot p_i=-\dfrac{\partial H}{\partial q_i}=-\dfrac{q_i}{|q|^3}
  \end{cases}
  \qquad\text{对 } i=1,2,3 .
\]
记 $e$ 为指向某方向的单位向量。由牛顿第二定律，中心力正比于 $q$ 对时间 $t$ 的二阶导数，
于是对某常数 $k\in\R$，有如下关系
\[
  F=-k\dot p=k\cdot\frac{q}{|q|^3}=\frac{k}{|q|^2}\cdot\frac{q}{|q|}
  =\frac{k}{|q|^2}e .
\]
这表明中心力与距离的平方成反比。\hfill$//$
\end{wexam}
% [Tutor Note: 这里用到 $\ddot q=\dot p$（因 $\dot q=p$）。符号上须留意：由 (2.31) ，
% $\dot p=-q/|q|^3$，故 $F=-k\dot p=k\,q/|q|^3$；而 (2.31) 对应的是\emph{吸引}力
% $F=-q/|q|^3$，即上式中 $k=-1$。原书让 $k$ 取任意实常数，故结论「与距离平方成反比」
% 成立，但读者若取 $k>0$ 会得到排斥力，与势函数 $-1/|q|$ 的物理含义相反。]

接下来，为验证经典结果「$(T^*(\R^3\setminus\{0\}),\omega_{\mathrm{can}},H)$ 是可积
系统」，由定义 2.6.1 需要三个哈密顿函数 $H_1(=H),H_2,H_3$，使其微分线性无关且两两
Poisson 交换。求 $H_2$ 相对直接（参见例 2.7.1）。取 $H_2=q_1p_2-q_2p_1$，则
\begin{align*}
  \{H,H_2\}
  &=\sum_{i=1}^2\left(\frac{\partial H}{\partial p_i}\frac{\partial H_2}{\partial q_i}
    -\frac{\partial H}{\partial q_i}\frac{\partial H_2}{\partial p_i}\right)\\
  &=p_1p_2-\left(-\frac{q_1}{|q|^3}\right)\cdot(-q_2)+p_2(-p_1)
    -\left(-\frac{q_2}{|q|^3}\right)\cdot q_1=0 .
\end{align*}
注意 $H_2$ 表示角动量 $A(=q\times p)=(A_1,A_2,A_3)$ 的一个分量（见 (2.26)），更确切地说
$H_2=A_3$。这些分量满足如下关系
\[
  \{A_1,A_2\}=-A_3,\qquad
  \{A_2,A_3\}=-A_1,\qquad
  \{A_1,A_3\}=A_2 .
\]
因此，无论取 $H_3=A_1$ 还是 $H_3=A_2$，就可积性而言都不行。
% [Tutor Note: 易错预警——按 (2.8) 与 (2.5) 的约定，$\{F,G\}=\sum_i(\partial_{q_i}F
% \partial_{p_i}G-\partial_{p_i}F\partial_{q_i}G)$（即通常的标准括号），此时角动量满足
% $\{A_1,A_2\}=A_3$、$\{A_2,A_3\}=A_1$、$\{A_1,A_3\}=-A_2$；原文列出的三个等式整体与
% 之反号，等价于采用相反符号的括号约定。由于本段只需 $\{A_i,A_j\}\neq0$ 这一事实，
% 结论不受影响。同理，上面 $\{H,H_2\}$ 的和式上限应为 $3$（$i=3$ 项自动为零），
% 且其被求和式整体与标准括号反号，但结果仍为 $0$。详见 review_ch02.md。]
% >>> p052
事实证明，找出第三个哈密顿函数 $H_3$ 相当有挑战性。历史上，\textbf{Runge--Lenz 向量}
（Runge--Lenz vector）被发现并定义为
\begin{equation}
  V_{\mathrm{RL}}(q,p):=p\times A-\frac{q}{|q|}\in\R^3,
  \tag{2.32}
\end{equation}
其中 $p\times A$ 是 $p$ 与角动量 $A$（在 (2.26) 中等于 $q\times p$）的叉积。下面的练习
说明，借助 $V_{\mathrm{RL}}$ 可得到哈密顿系统 $(T^*(\R^3\setminus\{0\}),
\omega_{\mathrm{can}},H)$ 的可积性。

\begin{wexer}{2.6.1}
（[86] 中的引理 3.3.5）记
$V_{\mathrm{RL}}=\big((V_{\mathrm{RL}})_1,(V_{\mathrm{RL}})_2,(V_{\mathrm{RL}})_3\big)$。
验证 $\{H,(V_{\mathrm{RL}})_i\}=0$（$i=1,2,3$），其中 $H$ 由 (2.31) 定义。
\end{wexer}
% [Tutor Note: 由练习 2.6.1，$V_{\mathrm{RL}}$ 的三个分量都是守恒量，但它们并不两两
% Poisson 交换：在 $SO(3)$ 作用下 $V_{\mathrm{RL}}$ 按向量变换，故
% $\{A_i,(V_{\mathrm{RL}})_j\}=\varepsilon_{ijk}(V_{\mathrm{RL}})_k$（符号随括号约定而
% 整体改变）。取 $H_1=H$、$H_2=A_3$、$H_3=(V_{\mathrm{RL}})_3$ 便有
% $\{A_3,(V_{\mathrm{RL}})_3\}=0$，且三者微分在开稠密集上线性无关，故确实给出 Kepler
% 问题的可积系统。$H$、$A$、$V_{\mathrm{RL}}$ 之间的括号关系构成 $\mathfrak{so}(4)$
% 型代数，这正是 Kepler 问题「隐藏对称性」的来源。]

下面是一个简单得多的例子，它与例 2.5.2 密切相关。

\begin{wexam}{2.6.2}
对 $(\C^n,\omega_{\mathrm{std}})$，考虑如下哈密顿函数族
\begin{equation}
  H_i(z_1,\cdots,z_n)=\frac12|z_i|^2\qquad\text{对 } i=1,\cdots,n .
  \tag{2.33}
\end{equation}
则显然 $H_1,\cdots,H_n$ 构成一个可积系统（对由 $H=H_i$（任意 $i$）给出的哈密顿系统
而言）。

此外，还可考虑由 (2.33) 中 $H_i$ 的线性组合给出的哈密顿系统。例如取
$H(z_1,\cdots,z_n)=\sum\limits_{i=1}^nH_i$，其中 $H_i$ 由 (2.33) 定义。这样的 $H$ 仍可
建立可积系统 $(\C^n,\omega_{\mathrm{std}},H)$。在例 2.7.2 中我们将看到，(2.33) 的考虑
是相当自然的。

最后，对任意在 (2.27) 的标准 $\T^n$-作用下不变的开子集 $U\subset\C^n$，其限制
$(U,\omega_{\mathrm{std}}|_U)$ 也是辛流形，且 $H$ 在 $U$ 上有定义。上述论证同样成立，
从而得到可积系统 $(U,\omega_{\mathrm{std}}|_U,H)$。在 7.2 节我们将看到，这样的 $U$
（及其闭包）是辛几何中所研究的基本对象。\hfill$//$
\end{wexam}

下面是一个稍复杂些的可积系统例子，它推广了例 2.6.2。

\begin{wexam}{2.6.3}
回顾由定义 2.5.2，辛环面流形 $(M^{2n},\omega)$ 容许有效的哈密顿 $\T^n$-作用。对其矩
映射 $\mu=(\mu_1,\cdots,\mu_n):M\to\R^n$——其中 $\R^n$ 与 $\T^n$ 的 Lie 代数的对偶
等同——把每个 $\mu_i:M\to\R$ 看作 $(M,\omega)$ 上的哈密顿函数，则
$(M,\omega,\mu_i)$ 是可积系统（对任意 $1\le i\le n$）。事实上，对任意 $i,j$ 有
\[
  \{\mu_i,\mu_j\}=\omega(X_{\mu_j},X_{\mu_i})=0 .
\]
这里第一个等号来自 (2.9)；第二个等号（等于零）是复习定义 2.5.1 中矩映射的对偶式定义的
好练习，它最终归结为方程 (2.23)。因此，哈密顿函数组 $\mu_1,\cdots,\mu_n$ 有助于在
$(M,\omega)$ 上构成一个可积系统。\hfill$//$
\end{wexam}
% [Tutor Note: 这里 $X_{\mu_i}$ 指点 $\mu_i$ 的哈密顿向量场；由矩映射的定义
% $\dd\langle\mu,v\rangle=-\iota_{X_v}\omega$，取 $v$ 为 $\T^n$ 的第 $i$ 个基向量即得
% $X_{\mu_i}=X_{v_i}$（作用生成的向量场）。于是
% $\{\mu_i,\mu_j\}=\omega(X_{v_j},X_{v_i})$，而由 (2.23) 对环面作用的可交换性知其为
% $0$——这正是 (2.23) 的另一读法。]
% >>> p053
也有许多不可积的哈密顿系统。例如经典力学中的平面三体问题一般是不可积的（见 [149]
以及其定理 1.1 后方讨论中关于不可积性的相关史实）。

\begin{wexam}{2.6.4}
（平面三体问题）此平面问题研究在平面内运动的三个物体 $X_1,X_2,X_3$，其中每个 $X_i$
在 $T^*\R^2$ 中的运动由坐标 $(q_i,p_i)\in T^*\R^2$ 描述。借助哈密顿力学的表述，
$(T^*\R^2\times T^*\R^2\times T^*\R^2,\omega_{\mathrm{prod}})$ 上的哈密顿函数 $H$
定义为
\[
  H=\sum_{i=1}^3\frac{1}{2m_i}|p_i|^2
  -\sum_{1\le i<j\le3}\frac{m_im_j}{|q_i-q_j|},
\]
其中 $m_i$ 表示 $X_i$ 的质量（$1\le i\le3$）。

Poincaré 发现，一般说来 $H$ 的哈密顿方程难以求解，这些 ODE 不存在一般的解析解。然而，
若考虑限制性三体问题——例如三个物体中包含两个绕其质心运动的大质量物体和一个较小的
第三质量——则哈密顿方程可以大大简化。更多细节见 [86]。\hfill$//$
\end{wexam}
